\documentclass[twocolumn]{aastex631}

\shorttitle{Gaia DR3 asteroid rotation periods}
\shortauthors{Tropynin}

\begin{document}

\title{Rotation Periods of 41{,}842 Asteroids from Gaia DR3 Sparse Photometry, with Calibrated Reliability and a Null Test}

\author{Oleg Tropynin}
\affiliation{Independent researcher, Kazakhstan}

\begin{abstract}
I derive rotation periods for every numbered asteroid with at least 25 Gaia DR3 transits (42{,}195
objects; 41{,}842 solved). I use a periodogram with season offsets, a linear phase term and a
second-order Fourier series, with the rotational phase corrected for the phase-angle bisector, and Gaia-centric geometry from numerical orbit integration. Asteroids with
a known period are split once into a tuning half and a test half. Every choice, including the period-doubling
rule and a logistic model for the probability that a period is correct, is fixed on the tuning half or on a
null test in which the residuals are permuted within seasons. On the 4{,}834 test asteroids, confident (tier A)
periods are correct within 3\% in 87.8\% of cases (1540/1753), and the probabilities are well calibrated.
Without the null test, a threshold chosen for 80\% precision on known asteroids would have let 19.8\% of signal-free
asteroids through. On known asteroids that resemble the unknown population, tier A is 87\% correct (242/278).
Agreement with three other Gaia DR3 analyses is 98--99\%, and with TESS (Vavilov \& Carry) 90\%. I list 1{,}284 tier-A periods for asteroids with no
published period in LCDB, DAMIT, about 180 VizieR tables or the Minor Planet Bulletin. The pipeline processes
DR3 in about twenty minutes on ten cores and is intended to run on Gaia DR4 on release day.
\end{abstract}

\keywords{Asteroid rotation (2211) --- Light curves (918) --- Main belt asteroids (2036) --- Photometry (1234)}

\section{Introduction}

Gaia DR3 gave sparse photometry for about 157{,}000 numbered asteroids \citep{Tanga2023}. Three groups
have derived rotation periods or shapes from it. \citet{Durech2023} used convex inversion,
\citet{Cellino2024} used a genetic search and \citet{MacLennan2026} used a period search with
phase-curve modelling. All three use the same data, and a large fraction of the asteroids with enough
transits remains without a published period. Two things are needed before such periods can be used:
an honest estimate of how often a period is wrong, and the same estimate for the asteroids that are
\emph{not} already in other catalogues. This paper provides both and makes the pipeline fast enough
to rerun on Gaia DR4 (expected in December 2026, with about twice the time baseline).
A pilot on 600 asteroids was reported separately \citep{Tropynin2026a}.

\section{Data and geometry}

All 20 bulk files of \texttt{sso\_observation} (23.3\,M CCD rows) were averaged per
\texttt{transit\_id}, giving 3{,}069{,}170 transits of 156{,}789 numbered asteroids. The median
asteroid has 17 transits; 42{,}195 have at least 25 and are analysed here. Each transit takes the largest
relative flux error of its CCDs, with a 0.03\,mag floor.

Heliocentric positions come from osculating elements (JPL, epoch 2026 June) integrated back to
2014--2017 with a fourth-order Runge--Kutta scheme. The step is one day, and the integration includes
the Sun and the four giant planets, using JPL Horizons vectors. Gaia's position is given in the data.
Light time is iterated. Against JPL Horizons on 19 asteroids, the reduced magnitude agrees to 0.0007\,mag,
the observer distance to 0.0005\,au and the phase angle to 0.009$^\circ$. A two-body solution was
rejected because its errors reach 0.13\,mag.

\section{Period search}

The reduced magnitude $G-5\log_{10}(r\Delta)$ is fitted by weighted least squares. The model has one
constant per observing season (seasons are separated by gaps longer than 60\,d), a term linear in phase
angle and a second-order Fourier series in rotational phase. Outliers beyond 5 median absolute
deviations from the no-rotation model are removed in two passes.

Over the 34-month baseline the rotational phase seen from Gaia drifts with the longitude
$\lambda_{\rm PAB}$ of the phase-angle bisector. The Fourier terms therefore use the phase
$2\pi(ft - s\,\lambda_{\rm PAB}/360^\circ)$, where $s=\pm1$ is the unknown sense of rotation.
Both senses are searched and the better fit is kept. On the test set below, this correction raised
the fraction of correct periods from 45.5\% to 52.6\%. It also raised the number of confident periods
from 1250 to 1753 at higher precision, and it reduced the median difference from the sidereal periods
of \citet{Durech2023} by a factor of four.

The search scans periods of 2--1000\,h on a grid with a step of $0.5/T$, where $T$ is the time span. The
30 deepest minima are then refined. A grid five times finer changed test precision by less than one
point.

Most lightcurves are double-peaked, so at the true period the second harmonic dominates. If the
first-harmonic amplitude $A_1$ exceeds $1.20A_2$, the doubled period is adopted. The factor 1.20 was
chosen on the tuning half. A BIC comparison of fourth-order models at $P$ and $2P$ was tried first and
rejected, because it favoured the half period.

\section{Calibration, null test and tiers}

The truth set is LCDB V4.0 \citep{Warner2009}, quality $U\ge2$: 9{,}613 of the solved asteroids. A
period counts as correct if it is within 3\% of the LCDB value. The asteroids are split by a hash of
their number into TUNE (4{,}779) and TEST (4{,}834). On TUNE I fitted a logistic model of the probability
of a correct period. Its inputs are the $\chi^2$ gain over the no-rotation model, the separation from the
next non-alias minimum, the number of transits, the amplitude, the reduced $\chi^2$, the period and
$A_1/A_2$. On TEST the predicted probabilities match the observed fractions (Figure~\ref{fig:rel}).

A threshold that gives 80\% precision on TUNE (0.384) is not safe for unknown asteroids. To test it, I
permuted the residuals of the no-rotation model within each season for 4{,}000 random unknown asteroids.
This keeps the sampling, the errors and the season levels but destroys rotation. At 0.384, 19.8\% of these
signal-free asteroids still pass. At 0.70 the rate falls to 0.85\% (Figure~\ref{fig:thr}). Tier~A is
therefore defined as probability $\ge0.70$ and period $<100$\,h. Tier~B covers 0.384--0.70. Best periods
of 100--1000\,h are flagged as slow candidates and never treated as confident, because only about 40\% of
them are correct.

\section{Results}

Of 41{,}842 solutions, 7{,}911 are tier~A and 6{,}102 are tier~B; 1{,}162 are flagged as slow candidates.
On TEST the results are:
\begin{itemize}
\item tier A: 1540 of 1753 correct (87.8\%), with 89\%, 87\% and 84\% in the 2--10, 10--30 and 30--100\,h bands;
\item the 213 tier-A errors: 200 other periods, 11 at $2P$ and 2 at $P/2$;
\item tier B: 62\% correct.
\end{itemize}

The comparison with LCDB is shown in Figure~\ref{fig:lcdb}. Correct periods agree with LCDB to a median of
$2.3\times10^{-4}$ and with the sidereal periods of \citet{Durech2023} to $3.9\times10^{-5}$. The curvature
error of the $\chi^2$ minimum underestimates the latter scatter. I therefore add $5.5\times10^{-5}P$ in
quadrature, which gives a median $|\Delta P|/\sigma$ of about 0.67.

\emph{Transfer to unknown asteroids.} The TEST asteroids are brighter and better studied than the
asteroids without a period. The closest known analogue is the set of TEST asteroids that appear in none of
the six catalogues of Table~\ref{tab:lit}. Tier~A is 87\% correct on them (242/278; 95\% interval
83--91\%). I adopt this figure for the new periods.

\begin{deluxetable}{lccc}
\tablecaption{Agreement of tier-A periods with published catalogues (within 3\%).\label{tab:lit}}
\tablehead{\colhead{Catalogue} & \colhead{Data} & \colhead{$n$} & \colhead{Agree}}
\startdata
\citet{Durech2023} & Gaia DR3 & 4287 & 98.9\% \\
\citet{Cellino2024} & Gaia DR3 & 3402 & 98.1\% \\
\citet{MacLennan2026} & Gaia DR3 & 3178 & 98.0\% \\
\citet{Vavilov2025} & TESS & 926 & 89.6\% \\
\citet{Erasmus2020} & ATLAS & 334 & 66.5\% \\
\citet{McNeill2023} & TESS & 1078 & 52.6\% \\
\enddata
\tablecomments{Rows that use Gaia DR3 test consistency, not accuracy. For TEST asteroids that are also in
LCDB, this work and the catalogue agree with LCDB in 435 vs 419 of 487 cases (Vavilov \& Carry), 133 vs 135 of
175 (Erasmus et al.), and 411 vs 273 of 447 (McNeill et al.).}
\end{deluxetable}

\emph{New periods.} 1{,}400 tier-A asteroids have no period in LCDB or in the six catalogues. 116 of them do
have a period in DAMIT \citep{Durech2010}, in about 180 VizieR tables with rotation periods (5.0\,M rows) or
in the Minor Planet Bulletin, volumes 50--53. Where such a value exists, this work agrees with it in 63 of
89 cases. The remaining 1{,}284 (977 at 2--10\,h, 201 at 10--30\,h, 106 at 30--100\,h) are listed as new, with
an expected precision of about 87\%. Behrend's Geneva page and ZTF catalogues outside VizieR were not
searched.

\emph{Pilot candidates.} Of the 42 candidates of \citet{Tropynin2026a}:
\begin{itemize}
\item 35 agree with this work;
\item 4 flagged ``low'' are doubled here;
\item for 3 flagged ``high'' (39932, 109412, 13686), the extended 1000\,h search prefers slow solutions, none of them tier~A. These three should be treated as ambiguous.
\end{itemize}

\section{Limitations and plans for Gaia DR4}

The $A_1/A_2$ rule is a heuristic. A model-based comparison of $P$ and $2P$ (ellipsoid or convex fit)
would be more robust. The PAB correction assumes a fixed sense of rotation and no pole information; a coarse pole search would refine it.
Tumbling and slow rotators fall outside a single-period model. For DR4, all thresholds, the null test and
the matched-control test will be fixed in advance, and injection--recovery tests on the real DR4
sampling will be added, so that results can be released within a day without re-tuning.

\begin{figure}
\plotone{fig1_lcdb.pdf}
\caption{Period from this work against the LCDB period for TEST asteroids. Tier A is shown in blue and
tier B in orange. The lines mark $P$, $2P$ and $P/2$.\label{fig:lcdb}}
\end{figure}

\begin{figure}
\plotone{fig2_reliability.pdf}
\caption{Reliability diagram on TEST, all solutions. Bars are 95\% binomial intervals; numbers give the
count per bin.\label{fig:rel}}
\end{figure}

\begin{figure*}
\plotone{fig3_threshold.pdf}
\caption{Left: fraction of signal-free (permuted) asteroids that pass as a function of the probability
threshold. Right: precision on TEST. Dashed: the TUNE 80\%-precision threshold. Solid: the tier-A
threshold.\label{fig:thr}}
\end{figure*}

\section{Data availability}

The catalogue (41{,}842 rows), the list of 1{,}284 new periods, the validation tables and the code will be archived on Zenodo [DOI to be assigned].

\begin{acknowledgments}
This work used data from the ESA mission Gaia, processed by the Gaia DPAC; JPL Horizons and the JPL
Small-Body Database; LCDB; and VizieR (CDS, Strasbourg). The analysis was designed, coded and checked with
the help of an AI assistant (Claude, Anthropic) under the author's direction. Two AI-based referee passes
were used to stress-test the validation before submission.
\end{acknowledgments}

\begin{thebibliography}{}
\bibitem[Cellino et al.(2024)]{Cellino2024} Cellino, A., et al. 2024, A\&A, 687, A277
\bibitem[\v{D}urech \& Hanu\v{s}(2023)]{Durech2023} \v{D}urech, J., \& Hanu\v{s}, J. 2023, A\&A, 675, A24
\bibitem[\v{D}urech et al.(2010)]{Durech2010} \v{D}urech, J., Sidorin, V., \& Kaasalainen, M. 2010, A\&A, 513, A46
\bibitem[Erasmus et al.(2020)]{Erasmus2020} Erasmus, N., et al. 2020, ApJS, 247, 13
\bibitem[MacLennan et al.(2026)]{MacLennan2026} MacLennan, E. M., et al. 2026, A\&A, 707, A131
\bibitem[McNeill et al.(2023)]{McNeill2023} McNeill, A., et al. 2023, AJ, 166, 152
\bibitem[Tanga et al.(2023)]{Tanga2023} Gaia Collaboration, Tanga, P., et al. 2023, A\&A, 674, A12
\bibitem[Tropynin(2026)]{Tropynin2026a} Tropynin, O. 2026, Candidate rotation periods for 42 main-belt asteroids from Gaia DR3 sparse photometry, Zenodo, doi:\dataset[10.5281/zenodo.23624115]{https://doi.org/10.5281/zenodo.23624115}
\bibitem[Vavilov \& Carry(2025)]{Vavilov2025} Vavilov, D., \& Carry, B. 2025, A\&A, 693, A66
\bibitem[Warner et al.(2009)]{Warner2009} Warner, B.~D., Harris, A.~W., \& Pravec, P. 2009, Icarus, 202, 134
\end{thebibliography}

\end{document}
